\subsection{Equilibrium Characterization}

We now characterize the equilibrium cutoff $x^*$ more formally. Define the posterior belief of agent $i$ as:
\begin{equation}
\hat{\theta}_i = \mathbb{E}[\theta \mid x_i, s].
\end{equation}

Let $\pi(\hat{\theta}_i)$ denote the perceived probability of crisis. The marginal agent at the cutoff $x^*$ must be indifferent between attacking and not attacking. This implies the indifference condition:
\begin{equation}
U^{attack}(\hat{\theta}^*) = U^{no\ attack}(\hat{\theta}^*).
\end{equation}

Using expected utility, this can be written as:
\begin{equation}
\pi(\hat{\theta}^*) \cdot B(\theta) = C,
\end{equation}

where $B(\theta)$ is the benefit of exiting early (avoiding crisis losses) and $C$ is the marginal cost of acquiring foreign currency.

\subsection{Comparative Statics}

We now study how equilibrium outcomes change with the precision of the public signal.

\paragraph{Effect on cutoff}

\textbf{Proposition 5.} An increase in public signal precision lowers the cutoff $x^*$.

\textit{Proof.}

Higher precision increases the weight on the public signal. Since the public signal is common across agents, beliefs become more aligned. This reduces uncertainty about others' actions and makes agents more willing to attack for a given private signal realization. Therefore, the cutoff $x^*$ decreases. \hfill $\square$

\paragraph{Effect on aggregate attack elasticity}

\textbf{Proposition 6.} The slope of the aggregate attack function $R(\theta)$ increases with public signal precision.

\textit{Proof.}

The derivative of $R(\theta)$ is:
\begin{equation}
\frac{dR}{d\theta} = \frac{1}{\sigma_x} \phi\left(\frac{x^* - \theta}{\sigma_x}\right),
\end{equation}

where $\phi(\cdot)$ is the standard normal density. A lower cutoff $x^*$ shifts mass toward the threshold region, increasing the slope locally. Thus, the system becomes more sensitive to fundamentals. \hfill $\square$

\subsection{Crisis Amplification}

\textbf{Proposition 7.} Endogenous increases in public signal precision amplify crisis probability.

\textit{Proof.}

Combining Propositions 2, 5, and 6, an increase in stablecoin usage raises signal precision, lowers the cutoff, and increases the responsiveness of aggregate attacks. This raises the probability that the crisis condition is satisfied for a given fundamental $\theta$. \hfill $\square$

\subsection{Welfare Decomposition}

We decompose welfare into two components:
\begin{equation}
W = W^{efficiency} - W^{fragility}.
\end{equation}

\paragraph{Efficiency term}
\begin{equation}
W^{efficiency} = \mathbb{E}[\log(c^N)] - \text{transaction costs}.
\end{equation}

\paragraph{Fragility term}
\begin{equation}
W^{fragility} = \mathbb{E}[\pi \cdot (\log(c^N) - \log(c^C))].
\end{equation}

\textbf{Proposition 8.} There exists a threshold $\theta^*$ such that stablecoins increase welfare for $\theta < \theta^*$ and decrease welfare for $\theta > \theta^*$.

\textit{Proof.}

At low $\theta$, crisis probability is small and the fragility term is negligible, so efficiency gains dominate. At high $\theta$, the increase in crisis probability raises expected losses, dominating efficiency gains. Continuity implies the existence of a threshold. \hfill $\square$

\subsection{Discussion}

These results formalize the central tradeoff in the model. Stablecoins improve information and reduce costs, but the endogenous increase in signal precision strengthens coordination and amplifies fragility. The equilibrium effects are therefore inherently non-linear and depend critically on the level of underlying misalignment.
